% ============================================================
% PRM393 - Section 1: Introduction (Bilingual)
% ============================================================
\section{\bi{Mở đầu}{Introduction}}
\label{sec:intro}

\bi{
Trong kỷ nguyên phân phối nội dung số hiện đại, ứng dụng di động đã trở
thành điểm tiếp xúc chính yếu giữa người dùng và các dịch vụ truyền thông
trực tuyến. Khác với môi trường máy tính để bàn (web desktop) có đường
truyền mạng ổn định và nguồn năng lượng dồi dào, môi trường di động đặt ra
những thách thức kỹ thuật nghiêm ngặt: độ trễ mạng biến động liên tục, kết
nối thường xuyên bị gián đoạn khi chuyển vùng (cell tower handover / Wi-Fi
to cellular), hạn chế về bộ nhớ RAM và thời lượng pin thiết bị.

Để giải quyết toàn diện các thách thức trên, bài báo này trình bày thiết
kế và hiện thực ứng dụng di động \textbf{OmniCast Mobile} --- một ứng dụng
phát thanh truyền hình, truyền phát trực tiếp (LiveStream) và phân phối nội dung số đa nền tảng được nghiên cứu và phát triển tại \textbf{FPTU Quy Nhon AI Campus --- FU Cahata}
cho môn học \textbf{PRM393 --- Mobile Applications Development} dưới sự hướng dẫn chuyên môn của giảng viên (\emph{Mentor}) \textbf{Thầy Huỳnh Vang}. Ứng dụng
được xây dựng trên nền tảng Flutter 3.x với ngôn ngữ Dart (53.922 dòng mã Dart), áp dụng kiến trúc
phân tầng Clean Architecture 4 lớp kết hợp mô hình quản lý trạng thái BLoC (Business Logic
Component) điều phối 31 màn hình và 22 thành phần giao diện chuyên dụng.
}{
In the modern era of digital content delivery, mobile applications have
become the primary touchpoint connecting viewers with broadcast media
services. Unlike desktop environments with steady broadband connectivity
and continuous power supplies, mobile devices present stringent technical
challenges: highly fluctuating network latency, frequent network drops during
roaming transitions (cell tower handovers or Wi-Fi to cellular switches),
constrained memory footprints, and battery budget limits.

To comprehensively address these demands, this paper presents the
architectural design and implementation of \textbf{OmniCast Mobile} --- an
enterprise cross-platform digital broadcast, LiveStream and multimedia streaming client
developed at \textbf{FPTU Quy Nhon AI Campus --- FU Cahata} for the course \textbf{PRM393 --- Mobile Applications Development} under the academic mentorship of \textbf{Mr. Huỳnh Vang}.
Built upon Flutter 3.x with the Dart language (53,922 Dart LOC), the client adopts a clean
four-tier layered architecture coupled with the Business Logic Component
(BLoC) state management pattern across 31 screens and 22 custom widgets.
}

\subsection{\bi{Động lực và bài toán kỹ thuật}{Motivation and Engineering Challenges}}

\bi{
Quá trình xây dựng ứng dụng di động OmniCast tập trung vào việc giải quyết năm
bài toán kỹ thuật cốt lõi:
}{
The engineering of OmniCast Mobile resolves five pivotal technical
challenges:
}

\begin{enumerate}[leftmargin=*,itemsep=2pt]
  \item \bi{\textbf{Khả năng hoạt động khi mất mạng (Offline-First):}
  Người dùng thường xuyên tra cứu lịch phát sóng (EPG) và danh sách theo
  dõi (Watchlist) trong điều kiện mất kết nối. Ứng dụng phải hỗ trợ lưu
  trữ cục bộ bằng cơ sở dữ liệu SQLite, xếp hàng các thao tác phát sinh khi
  ngoại tuyến và tự động phát lại (replay) đồng bộ lên máy chủ khi có mạng
  trở lại.}{%
  \textbf{Offline-First Resilience:} Viewers frequently consult electronic
  programme guides (EPG) and their personal watchlists in poor or nonexistent
  connectivity conditions. The application must store structured data in a
  local SQLite database, queue pending mutations while offline, and seamlessly
  replay them to the central cloud upon reconnection.}

  \item \bi{\textbf{Luồng phát đa phương tiện độ trễ thấp và tự động khám phá mạng LAN:}
  Hệ sinh thái OmniCast phát sóng qua nhiều giao thức: WebRTC / HLS thông
  qua MediaMTX, luồng trực tiếp Google Drive và luồng nhúng video ngoài. Trình phát \texttt{OmniPlayer}
  phải tự động quét và khám phá địa chỉ IP máy chủ MediaMTX trong mạng Wi-Fi cục bộ (\texttt{LanDiscoveryService}) và tự phục hồi
  khi luồng video bị đứt quãng.}{%
  \textbf{Low-Latency Adaptive Streaming and LAN Host Discovery:} The OmniCast broadcast pipeline
  emits ultra-low-latency WebRTC/HLS streams via MediaMTX, direct Google Drive streams, and external
  curated feeds. The embedded \texttt{OmniPlayer} must dynamically scan and discover
  local Wi-Fi MediaMTX hosts (\texttt{LanDiscoveryService}) and recover transparently during transient
  network drops.}

  \item \bi{\textbf{Phòng xem LiveStream tương tác hai chiều \& Bàn điều khiển Mobile Live Studio:}
  Người xem không chỉ thụ động tiếp nhận luồng phát mà còn tham gia sảnh chờ, gửi bình luận phân nhánh,
  thả biểu cảm hạt nổi có phản hồi xúc giác 0ms, nhận thông báo ghim từ đạo diễn và quyên góp Super Chat 6 mệnh giá. Đồng thời, phóng viên hiện trường cần một bàn điều khiển \texttt{LiveStudioScreen} 4 tab và ống kính máy quay \texttt{MobileStudioRecorderScreen} hỗ trợ đẩy sóng WebRTC WHIP trực tiếp từ điện thoại.}{%
  \textbf{Bi-Directional Interactive LiveStream \& Mobile Live Studio:} Viewers expect rich interactive
  rooms with lobby prompts, threaded live chat, 0ms haptic floating reactions, pinned host banners, and 6-tier Super Chat donations. Meanwhile, field reporters require a 4-tab \texttt{LiveStudioScreen} and hardware camera viewfinder (\texttt{MobileStudioRecorderScreen}) capable of real-time WebRTC WHIP publishing.}

  \item \bi{\textbf{Quản lý trạng thái phân tán và bảo mật phiên:}
  Việc điều hướng 31 màn hình với 12 module BLoC yêu cầu cơ chế kiểm soát
  trạng thái chặt chẽ, cùng với việc lưu trữ token JWT an toàn trong
  Keystore/Keychain để tự động gia hạn phiên làm việc ngầm mà không làm phiền
  người dùng.}{%
  \textbf{Distributed State Management and Session Security:}
  Navigating 31 screens powered by 12 BLoC modules requires predictable
  unidirectional data flows, paired with hardware-backed secure storage
  (Keystore/Keychain) for transparent JWT token rotation.}
\end{enumerate}

\subsection{\bi{Tổ chức nhóm nghiên cứu và Phân công nhiệm vụ}{Project Team Organization and Role Allocation}}

\bi{
Dự án được thực hiện bởi nhóm 4 sinh viên tại \textbf{FPTU Quy Nhon AI Campus --- FU Cahata} dưới sự hướng dẫn của giảng viên \textbf{Thầy Huỳnh Vang}, như trình bày tại Bảng~\ref{tab:team_prm}.
}{
The project was executed by a four-member student engineering team at \textbf{FPTU Quy Nhon AI Campus --- FU Cahata} under the supervision of Mentor \textbf{Mr. Huỳnh Vang}, as detailed in Table~\ref{tab:team_prm}.
}

\begin{table}[t]
\caption{\bi{Danh sách 4 thành viên nhóm dự án OmniCast (FPTU Quy Nhon AI Campus --- FU Cahata, Mentor: Thầy Huỳnh Vang).}{The 4 project team members of OmniCast (FPTU Quy Nhon AI Campus --- FU Cahata, Mentor: Mr. Huỳnh Vang).}}
\label{tab:team_prm}
\centering
\scriptsize
\begin{tabular}{@{}p{2.3cm}p{1.5cm}p{4.3cm}@{}}
\toprule
\textbf{\bi{Họ và tên}{Full Name}} & \textbf{\bi{Vai trò / MSSV}{Role / ID}} & \textbf{\bi{Phụ trách kỹ thuật chính}{Primary Engineering Ownership}} \\
\midrule
\textbf{Nguyễn Bình An} (\texttt{TomOutfit}) & \bi{Nhóm trưởng}{Team Leader} & \bi{Kiến trúc lõi NestJS 19 module, Bảo mật JWT \& RBAC 4 vai trò, Thuật toán chống xung đột lịch EPG, Direct Stream Engine, CI/CD \& Triển khai}{NestJS 19-module core architecture, JWT \& 4-role RBAC security, EPG conflict detection SQL, Direct Stream Engine, CI/CD \& Release} \\
\addlinespace[2pt]
\textbf{Trần Thanh Sơn} & \texttt{QE170080} & \bi{Kiến trúc vi dịch vụ gRPC Audit Logger (\texttt{audit\_logger.proto}), Cơ sở dữ liệu PostgreSQL Supabase (21 Prisma models), Telemetry \& Seed}{Standalone gRPC Audit Logger microservice (\texttt{audit\_logger.proto}), Supabase PostgreSQL (21 Prisma models), Telemetry \& Seeding} \\
\addlinespace[2pt]
\textbf{Nguyễn Vĩnh Phú} & \texttt{QE180136} & \bi{Trưởng nhóm Mobile Flutter (31 màn hình, 12 BLoC), Trình phát \texttt{OmniPlayer}, \texttt{MobileStudioRecorderScreen} (WebRTC WHIP), Offline SQLite}{Mobile Flutter Lead (31 screens, 12 BLoCs), \texttt{OmniPlayer} engine, \texttt{MobileStudioRecorderScreen} (WebRTC WHIP), Offline SQLite} \\
\addlinespace[2pt]
\textbf{Đồng Hoàng Nguyên} (\texttt{teng122}) & Full-Stack \& AI & \bi{Cổng Web Next.js 15 Cinema Dark (33 route), Cổng Socket.IO \texttt{/live}, Phòng Live 7 chủ đề, Bàn điều khiển Live Studio \& Module AI Curator}{Next.js 15 Cinema Dark Web Portal (33 routes), Socket.IO \texttt{/live} gateway, 7-theme Live Room, Live Studio Console \& AI Curator module} \\
\bottomrule
\end{tabular}
\end{table}

\subsection{\bi{Đóng góp chính của báo cáo}{Principal Contributions}}

\bi{
Báo cáo này phân tích chi tiết năm đóng góp công nghệ của ứng dụng:
}{
This paper details five principal technical contributions:
}

\begin{itemize}[leftmargin=*,itemsep=2pt]
  \item \bi{Kiến trúc Clean Architecture 4 tầng (Presentation, Logic/BLoC,
  Data, Core) với cơ chế Service Locator qua \texttt{get\_it}.}{%
  A 4-tier Clean Architecture (Presentation, Logic/BLoC, Data, Core)
  orchestrated through \texttt{get\_it} service locator.}
  \item \bi{Hệ sinh thái \texttt{OmniPlayer}, \texttt{LiveWatchScreen}, \texttt{LiveStudioScreen} (4 tab) và máy quay \texttt{MobileStudioRecorderScreen} (WebRTC WHIP, SMPTE Timecode, Stereo VU Meter, Teleprompter HUD).}{%
  A complete mobile live broadcast ecosystem comprising \texttt{OmniPlayer}, \texttt{LiveWatchScreen}, the 4-tab \texttt{LiveStudioScreen} and the \texttt{MobileStudioRecorderScreen} field viewfinder (WebRTC WHIP, SMPTE timecode, Stereo VU meter, Teleprompter HUD).}
  \item \bi{Mô hình quản lý trạng thái luồng (Stream-based BLoC) với các
  concurrency transformer như \texttt{droppable} và \texttt{restartable}.}{%
  A stream-based reactive state architecture employing custom BLoC
  concurrency transformers such as \texttt{droppable} and \texttt{restartable}.}
  \item \bi{Động cơ ngoại tuyến SQLite với hàng đợi \texttt{pending\_ops} và cơ chế tự
  động giải quyết xung đột dữ liệu.}{%
  An offline-first SQLite synchronization engine featuring an asynchronous
  mutation queue (\texttt{pending\_ops}) and conflict resolution heuristics.}
  \item \bi{Bộ kiểm thử tự động toàn diện gồm 91 test đạt tỷ lệ vượt qua
  100\% trên quy trình CI/CD GitHub Actions.}{%
  A test suite of 91 automated unit and widget tests achieving a 100\%
  pass rate in a GitHub Actions continuous integration matrix.}
\end{itemize}
